\section{Affine morphisms}
\label{section:II.1}

Most of the results of this section are the ``global''
counterparts of those of Chapter~I, \S~1; they are
therefore not essentially new, and simply provide a convenient
language for what follows.

\setcounter{subsection}{0}
\subsection{$S$-preschemes and $\mathcal{O}_S$-algebras}
\label{subsection:II.1.1}

\begin{env}[1.1.1]
\label{II.1.1.1}
Let $S$ be a prescheme, $X$ an $S$-prescheme, and $f:X\to S$ its structural
morphism. We know (\hyperref[0.4.2.4]{0, 4.2.4}) that the direct image
$f_*(\mathcal{O}_X)$ is an $\mathcal{O}_S$-algebra, which we
\oldpage[II]{6}
will denote by $\mathcal{A}(X)$ when no confusion can result; if $U$
is an open subset of $S$, we have
\[
  \mathcal{A}(f^{-1}(U))=\mathcal{A}(X)|U.
\]
Similarly, for every $\mathcal{O}_X$-module $\mathcal{F}$ (resp. every
$\mathcal{O}_X$-algebra $\mathcal{B}$), we will write
$\mathcal{A}(\mathcal{F})$ (resp. $\mathcal{A}(\mathcal{B})$) for the direct
image $f_*(\mathcal{F})$ (resp. $f_*(\mathcal{B})$), which is an
$\mathcal{A}(X)$-module (resp. an $\mathcal{A}(X)$-algebra), and not only
an $\mathcal{O}_S$-module (resp. $\mathcal{O}_S$-algebra).
\end{env}

\begin{env}[1.1.2]
\label{II.1.1.2}
Let $Y$ be a second $S$-prescheme, $g:Y\to S$ its structural morphism,
and $h:X\to Y$ an $S$-morphism; we therefore have the commutative diagram
\[
  \xymatrix{
    X\ar[rr]^h\ar[rd]_f && Y\ar[ld]^g\\
    & S
  }
  \tag{1.1.2.1}\label{II.1.1.2.1}
\]

By definition we have $h=(\psi,\theta)$, where
$\theta:\mathcal{O}_Y\to h_*(\mathcal{O}_X)=\psi_*(\mathcal{O}_X)$ is a
homomorphism of sheaves of rings; from it we deduce
(\hyperref[0.4.2.2]{0, 4.2.2}) a homomorphism of
$\mathcal{O}_S$-algebras
\[
  g_*(\theta):g_*(\mathcal{O}_Y)\to
  g_*(h_*(\mathcal{O}_X))=f_*(\mathcal{O}_X),
\]
in other words, a homomorphism of $\mathcal{O}_S$-algebras
$\mathcal{A}(Y)\to\mathcal{A}(X)$, which we will also denote by
$\mathcal{A}(h)$. If $h':Y\to Z$ is a second $S$-morphism, it is immediate
that $\mathcal{A}(h'\circ h)=\mathcal{A}(h)\circ\mathcal{A}(h')$. We have
therefore defined a \emph{contravariant functor} $\mathcal{A}(X)$ on the
category of $S$-preschemes, with values in the category of
$\mathcal{O}_S$-algebras.

Now let $\mathcal{F}$ be an $\mathcal{O}_X$-module, $\mathcal{G}$ an
$\mathcal{O}_Y$-module, and $u:\mathcal{G}\to\mathcal{F}$ an $h$-morphism,
that is (\hyperref[0.4.4.1]{0, 4.4.1}), a homomorphism of
$\mathcal{O}_Y$-modules $\mathcal{G}\to h_*(\mathcal{F})$. Then
\[
  g_*(u):g_*(\mathcal{G})\to g_*(h_*(\mathcal{F}))=f_*(\mathcal{F})
\]
is a homomorphism $\mathcal{A}(\mathcal{G})\to
\mathcal{A}(\mathcal{F})$ of $\mathcal{O}_S$-modules, which we will denote by
$\mathcal{A}(u)$; moreover, the pair
$(\mathcal{A}(h),\mathcal{A}(u))$ constitutes a \emph{di-homomorphism} from
the $\mathcal{A}(Y)$-module $\mathcal{A}(\mathcal{G})$ to the
$\mathcal{A}(X)$-module $\mathcal{A}(\mathcal{F})$.
\end{env}

\begin{env}[1.1.3]
\label{II.1.1.3}
The prescheme $S$ being fixed, we can regard the pairs
$(X,\mathcal{F})$, where $X$ is an $S$-prescheme and $\mathcal{F}$ an
$\mathcal{O}_X$-module, as forming a \emph{category}, by defining a
\emph{morphism} $(X,\mathcal{F})\to(Y,\mathcal{G})$ as a pair $(h,u)$,
where $h:X\to Y$ is an $S$-morphism and $u:\mathcal{G}\to\mathcal{F}$ an
$h$-morphism. We can then say that
$(\mathcal{A}(X),\mathcal{A}(\mathcal{F}))$ is a \emph{contravariant
functor} with values in the category whose objects are the pairs
consisting of an $\mathcal{O}_S$-algebra and a module over this algebra, and
whose morphisms are the di-homomorphisms.
\end{env}

\subsection{Preschemes affine over a prescheme}
\label{subsection:II.1.2}

\begin{definition}[1.2.1]
\label{II.1.2.1}
Let $X$ be an $S$-prescheme and $f:X\to S$ its structural morphism. We say
that $X$ is \emph{affine over $S$} if there exists a covering
$(S_\alpha)$ of $S$ by affine open subsets such that, for every $\alpha$, the
prescheme induced by $X$ on the open set $f^{-1}(S_\alpha)$ is affine.
\end{definition}

\begin{example}[1.2.2]
\label{II.1.2.2}
Every closed sub-prescheme of $S$ is an $S$-prescheme affine over $S$
(\hyperref[I.4.2.3]{I, 4.2.3} and
\hyperref[I.4.2.4]{4.2.4}).
\end{example}

\begin{remark}[1.2.3]
\label{II.1.2.3}
A prescheme $X$ affine over $S$ is not necessarily an affine scheme,
as the example $X=S$ shows (\hyperref[II.1.2.2]{1.2.2}).
On the other hand, if an affine scheme $X$ is an $S$-prescheme, $X$ is not
necessarily affine over
\oldpage[II]{7}
$S$ (see (\hyperref[II.1.3.3]{1.3.3})). However, let us recall that if
$S$ is a \emph{scheme}, every $S$-prescheme which is an affine scheme is
affine over $S$ (\hyperref[I.5.5.10]{I, 5.5.10}).
\end{remark}

\begin{proposition}[1.2.4]
\label{II.1.2.4}
Every $S$-prescheme which is affine over $S$ is separated over
$S$ (in other words, is an $S$-scheme).
\end{proposition}

This follows immediately from (\hyperref[I.5.5.5]{I, 5.5.5}) and
(\hyperref[I.5.5.8]{I, 5.5.8}).

\begin{proposition}[1.2.5]
\label{II.1.2.5}
Let $X$ be an $S$-scheme affine over $S$, and $f:X\to S$ the structural
morphism. For every open subset $U\subset S$, $f^{-1}(U)$ is affine over
$U$.
\end{proposition}

By virtue of Definition (\hyperref[II.1.2.1]{1.2.1}), we are reduced to the
case where $S=\operatorname{Spec}(A)$ and $X=\operatorname{Spec}(B)$ are
affine, and then $f=({}^a\varphi,\widetilde{\varphi})$, where
$\varphi:A\to B$ is a homomorphism. Since the $D(g)$, where $g\in A$, form a
basis of $S$, we are reduced to the case where $U=D(g)$; but we then know
that $f^{-1}(U)=D(\varphi(g))$ (\hyperref[I.1.2.2.2]{I, 1.2.2.2}), whence
the proposition.

\begin{proposition}[1.2.6]
\label{II.1.2.6}
Let $X$ be an $S$-scheme affine over $S$, and $f:X\to S$ the structural
morphism. For every quasi-coherent $\mathcal{O}_X$-module $\mathcal{F}$,
$f_*(\mathcal{F})$ is a quasi-coherent $\mathcal{O}_S$-module.
\end{proposition}

Taking into account (\hyperref[II.1.2.4]{1.2.4}), this follows from
(\hyperref[I.9.2.2]{I, 9.2.2, a)}).

In particular, the $\mathcal{O}_S$-algebra
$\mathcal{A}(X)=f_*(\mathcal{O}_X)$ is \emph{quasi-coherent}.

\begin{proposition}[1.2.7]
\label{II.1.2.7}
Let $X$ be an $S$-scheme affine over $S$. For every $S$-prescheme
$Y$, the map $h\mapsto\mathcal{A}(h)$ from the set
$\operatorname{Hom}_S(Y,X)$ to the set
$\operatorname{Hom}(\mathcal{A}(X),\mathcal{A}(Y))$
(\hyperref[II.1.1.2]{1.1.2}) is bijective.
\end{proposition}

Let $f:X\to S$, $g:Y\to S$ be the structural morphisms. Suppose first that
$S=\operatorname{Spec}(A)$ and $X=\operatorname{Spec}(B)$ are affine; we
must prove that for every homomorphism
$\omega:f_*(\mathcal{O}_X)\to g_*(\mathcal{O}_Y)$ of
$\mathcal{O}_S$-algebras, there exists one and only one $S$-morphism
$h:Y\to X$ such that $\mathcal{A}(h)=\omega$. By definition, for every open
subset $U\subset
S$, $\omega$ defines a homomorphism
\[
  \omega_U=\Gamma(U,\omega):
  \Gamma(f^{-1}(U),\mathcal{O}_X)\to
  \Gamma(g^{-1}(U),\mathcal{O}_Y)
\]
of $\Gamma(U,\mathcal{O}_S)$-algebras. In particular, for $U=S$, this
gives a homomorphism
$\varphi:\Gamma(X,\mathcal{O}_X)\to\Gamma(Y,\mathcal{O}_Y)$ of
$\Gamma(S,\mathcal{O}_S)$-algebras, to which there corresponds a
well-determined $S$-morphism $h:Y\to X$, since $X$ is affine
(\hyperref[I.2.2.4]{I, 2.2.4}). It remains to prove that
$\mathcal{A}(h)=\omega$, in other words, that for every open subset $U$ in a
basis of $S$, $\omega_U$ coincides with the algebra homomorphism $\varphi_U$
corresponding to the $S$-morphism $g^{-1}(U)\to f^{-1}(U)$ that is the
restriction of $h$.
We may restrict to the case where $U=D(\lambda)$, with $\lambda\in S$; then,
if $f=({}^a\rho,\widetilde{\rho})$, where $\rho:A\to B$ is a ring
homomorphism, we have $f^{-1}(U)=D(\mu)$, where $\mu=\rho(\lambda)$, and
$\Gamma(f^{-1}(U),\mathcal{O}_X)$ is the ring of fractions $B_\mu$; now,
the diagram
\[
  \xymatrix{
    B\ar[r]^\varphi\ar[d] & \Gamma(Y,\mathcal{O}_Y)\ar[d]\\
    B_\mu\ar[r]_{\varphi_U} & \Gamma(g^{-1}(U),\mathcal{O}_Y)
  }
\]
is commutative, and so is the analogous diagram in which $\varphi_U$
is replaced by $\omega_U$; the equality $\varphi_U=\omega_U$ then follows
from the universal property of rings of fractions
(\hyperref[0.1.2.4]{0, 1.2.4}).

Let us pass to the general case, and let $(S_\alpha)$ be a covering of $S$
by affine open subsets
\oldpage[II]{8}
such that the $f^{-1}(S_\alpha)$ are affine. Then every homomorphism
$\omega:\mathcal{A}(X)\to\mathcal{A}(Y)$ of $\mathcal{O}_S$-algebras gives,
by restriction, a family of homomorphisms
\[
  \omega_\alpha:
  \mathcal{A}(f^{-1}(S_\alpha))\to
  \mathcal{A}(g^{-1}(S_\alpha))
\]
of $\mathcal{O}_{S_\alpha}$-algebras, hence a family of
$S_\alpha$-morphisms
$h_\alpha:g^{-1}(S_\alpha)\to f^{-1}(S_\alpha)$ by the above.
Everything comes down to seeing that for every affine open subset $U$ in a
basis of $S_\alpha\cap S_\beta$, the restrictions of $h_\alpha$ and of
$h_\beta$ to $g^{-1}(U)$ coincide, which is evident since, by virtue of the
above, these restrictions both correspond to the homomorphism
$\mathcal{A}(X)|U\to\mathcal{A}(Y)|U$ that is the restriction of $\omega$.

\begin{corollary}[1.2.8]
\label{II.1.2.8}
Let $X$, $Y$ be two $S$-schemes which are affine over $S$. For an
$S$-morphism $h:Y\to X$ to be an isomorphism, it is necessary and sufficient
that $\mathcal{A}(h):\mathcal{A}(X)\to\mathcal{A}(Y)$ be an isomorphism.
\end{corollary}

This follows immediately from (\hyperref[II.1.2.7]{1.2.7}) and from the
functorial character of $\mathcal{A}(X)$.

\subsection{Prescheme affine over $S$ associated to an
$\mathcal{O}_S$-algebra}
\label{subsection:II.1.3}

\begin{proposition}[1.3.1]
\label{II.1.3.1}
Let $S$ be a prescheme. For every quasi-coherent $\mathcal{O}_S$-algebra
$\mathcal{B}$, there exists a prescheme $X$ affine over $S$, determined up
to a unique $S$-isomorphism, such that $\mathcal{A}(X)=\mathcal{B}$.
\end{proposition}

The uniqueness follows from (\hyperref[II.1.2.8]{1.2.8}); let us prove
the existence of $X$. For every affine open subset $U\subset S$, let $X_U$ be
the prescheme $\operatorname{Spec}(\Gamma(U,\mathcal{B}))$; since
$\Gamma(U,\mathcal{B})$ is a $\Gamma(U,\mathcal{O}_S)$-algebra, $X_U$
is an $S$-prescheme (\hyperref[I.1.6.1]{I, 1.6.1}). Moreover, since
$\mathcal{B}$ is quasi-coherent, the $\mathcal{O}_S$-algebra
$\mathcal{A}(X_U)$ is canonically identified with $\mathcal{B}|U$
(\hyperref[I.1.3.7]{I, 1.3.7},
\hyperref[I.1.3.13]{1.3.13} and
\hyperref[I.1.6.3]{1.6.3}). Let $V$ be a second affine open subset of
$S$, and let $X_{U,V}$ be the prescheme induced by $X_U$ on
$f_U^{-1}(U\cap V)$, where $f_U$ denotes the structural morphism
$X_U\to S$; $X_{U,V}$ and
$X_{V,U}$ are affine over $U\cap V$
(\hyperref[II.1.2.5]{1.2.5}), and by definition $\mathcal{A}(X_{U,V})$
and $\mathcal{A}(X_{V,U})$ are canonically identified with
$\mathcal{B}|(U\cap V)$. There is therefore
(\hyperref[II.1.2.8]{1.2.8}) a canonical $S$-isomorphism
$\theta_{U,V}:X_{V,U}\to X_{U,V}$; moreover, if $W$ is a third
affine open subset of $S$, and if $\theta'_{U,V}$, $\theta'_{V,W}$,
$\theta'_{U,W}$ are the restrictions of $\theta_{U,V}$, $\theta_{V,W}$,
$\theta_{U,W}$ to the inverse images of $U\cap V\cap W$ in $X_V$,
$X_W$ and $X_W$ respectively under the structural morphisms, we have
$\theta'_{U,V}\circ\theta'_{V,W}=\theta'_{U,W}$. Consequently there exist a
prescheme $X$, a covering $(T_U)$ of $X$ by affine open subsets, and
for each $U$ an isomorphism $\varphi_U:X_U\to T_U$, such that
$\varphi_U$ maps $f_U^{-1}(U\cap V)$ onto $T_U\cap T_V$ and that we have
$\theta_{U,V}=\varphi_U^{-1}\circ\varphi_V$
(\hyperref[I.2.3.1]{I, 2.3.1}). The morphism
$g_U=f_U\circ\varphi_U^{-1}$ makes $T_U$ into an $S$-prescheme, and the
morphisms $g_U$ and $g_V$ coincide in $T_U\cap T_V$, so $X$ is an
$S$-prescheme. Moreover, it is clear by definition that $X$ is affine
over $S$ and that $\mathcal{A}(T_U)=\mathcal{B}|U$, hence
$\mathcal{A}(X)=\mathcal{B}$.

We say that the $S$-scheme $X$ thus defined is \emph{associated to the
$\mathcal{O}_S$-algebra $\mathcal{B}$}, or is the \emph{spectrum of
$\mathcal{B}$}, and we denote it by $\operatorname{Spec}(\mathcal{B})$.

\begin{corollary}[1.3.2]
\label{II.1.3.2}
Let $X$ be a prescheme affine over $S$, and $f:X\to S$ the structural
morphism. For every affine open subset $U\subset S$, the prescheme induced
on $f^{-1}(U)$ is an affine scheme with ring
$\Gamma(U,\mathcal{A}(X))$.
\end{corollary}

\oldpage[II]{9}
Since we can suppose $X$ associated to an $\mathcal{O}_S$-algebra, by virtue
of (\hyperref[II.1.2.6]{1.2.6}) and (\hyperref[II.1.3.1]{1.3.1}), the
corollary follows from the construction of $X$ described in
(\hyperref[II.1.3.1]{1.3.1}).

\begin{example}[1.3.3]
\label{II.1.3.3}
Let $S$ be the affine plane over a field $K$ in which the point $0$ has been
doubled (\hyperref[I.5.5.11]{I, 5.5.11}); with the notation of
(\hyperref[I.5.5.11]{I, 5.5.11}), $S$ is the union of two affine open
subsets $Y_1$, $Y_2$; if $f$ is the open immersion $Y_1\to S$,
$f^{-1}(Y_2)=Y_1\cap Y_2$ is not an affine open subset in $Y_1$
(\emph{loc. cit.}), so we have an example of an affine scheme which is not
affine over $S$.
\end{example}

\begin{corollary}[1.3.4]
\label{II.1.3.4}
Let $S$ be an affine scheme; for an $S$-prescheme $X$ to be affine over
$S$, it is necessary and sufficient that $X$ be an affine scheme.
\end{corollary}

\begin{corollary}[1.3.5]
\label{II.1.3.5}
Let $X$ be a prescheme affine over a prescheme $S$, and $Y$ an
$X$-prescheme. For $Y$ to be affine over $S$, it is necessary and
sufficient that $Y$ be affine over $X$.
\end{corollary}

We are immediately reduced to the case where $S$ is an affine scheme, and
then so is $X$ (\hyperref[II.1.3.4]{1.3.4}); the two conditions of
the statement then express that $Y$ is an affine scheme
(\hyperref[II.1.3.4]{1.3.4}).

\begin{env}[1.3.6]
\label{II.1.3.6}
Let $X$ be a prescheme affine over $S$. Defining a prescheme $Y$
affine \emph{over $X$} amounts, by virtue of
(\hyperref[II.1.3.5]{1.3.5}), to giving a prescheme $Y$ affine
\emph{over $S$} and an $S$-morphism $g:Y\to X$; in other words
(\hyperref[II.1.3.1]{1.3.1} and
\hyperref[II.1.2.7]{1.2.7}), it amounts to giving a quasi-coherent
$\mathcal{O}_S$-algebra $\mathcal{B}$ and a homomorphism
$\mathcal{A}(X)\to\mathcal{B}$ of $\mathcal{O}_S$-algebras (which we can
regard as defining on $\mathcal{B}$ a structure of
$\mathcal{A}(X)$-algebra). If $f:X\to S$ is the structural morphism, we then
have $\mathcal{B}=f_*(g_*(\mathcal{O}_Y))$.
\end{env}

\begin{corollary}[1.3.7]
\label{II.1.3.7}
Let $X$ be a prescheme affine over $S$; for $X$ to be of finite type
over $S$, it is necessary and sufficient that the quasi-coherent
$\mathcal{O}_S$-algebra $\mathcal{A}(X)$ be of finite type
(\hyperref[I.9.6.2]{I, 9.6.2}).
\end{corollary}

By definition (\hyperref[I.9.6.2]{I, 9.6.2}), we are reduced to the case
where $S$ is affine; then $X$ is an affine scheme
(\hyperref[II.1.3.4]{1.3.4}), and if $S=\operatorname{Spec}(A)$,
$X=\operatorname{Spec}(B)$, $\mathcal{A}(X)$ is the
$\mathcal{O}_S$-algebra $\widetilde{B}$; since
$\Gamma(U,\widetilde{B})=B$, the corollary follows from
(\hyperref[I.9.6.2]{I, 9.6.2}) and
(\hyperref[I.6.3.3]{I, 6.3.3}).

\begin{corollary}[1.3.8]
\label{II.1.3.8}
Let $X$ be a prescheme affine over $S$; for $X$ to be reduced it is
necessary and sufficient that the quasi-coherent $\mathcal{O}_S$-algebra
$\mathcal{A}(X)$ be reduced
(\hyperref[0.4.1.4]{0, 4.1.4}).
\end{corollary}

Indeed, the question is evidently local on $S$; by virtue of
(\hyperref[II.1.3.2]{1.3.2}), the corollary follows from
(\hyperref[I.5.1.1]{I, 5.1.1}) and
(\hyperref[I.5.1.4]{I, 5.1.4}).

\subsection{Quasi-coherent sheaves on a prescheme affine over
$S$}
\label{subsection:II.1.4}

\begin{proposition}[1.4.1]
\label{II.1.4.1}
Let $X$ be a prescheme affine over $S$, $Y$ an $S$-prescheme, and
$\mathcal{F}$ (resp. $\mathcal{G}$) a quasi-coherent $\mathcal{O}_X$-module
(resp. $\mathcal{O}_Y$-module). Then the map
$(h,u)\mapsto(\mathcal{A}(h),\mathcal{A}(u))$ from the set of morphisms
$(Y,\mathcal{G})\to(X,\mathcal{F})$ to the set of di-homomorphisms
$(\mathcal{A}(X),\mathcal{A}(\mathcal{F}))\to
(\mathcal{A}(Y),\mathcal{A}(\mathcal{G}))$
(\hyperref[II.1.1.2]{1.1.2} and
\hyperref[II.1.1.3]{1.1.3}) is bijective.
\end{proposition}

The proof follows exactly the same lines as that of
(\hyperref[II.1.2.7]{1.2.7}), using
(\hyperref[I.2.2.5]{I, 2.2.5}) and
(\hyperref[I.2.2.4]{I, 2.2.4}), and the details are left to the
reader.

\begin{corollary}[1.4.2]
\label{II.1.4.2}
If, in addition to the hypotheses of (\hyperref[II.1.4.1]{1.4.1}), we
suppose $Y$ affine over $S$, then, for $(h,u)$ to be an isomorphism,
it is necessary and sufficient that $(\mathcal{A}(h),\mathcal{A}(u))$ be a
di-isomorphism.
\end{corollary}

\oldpage[II]{10}
\begin{proposition}[1.4.3]
\label{II.1.4.3}
For every pair $(\mathcal{B},\mathcal{M})$ consisting of a quasi-coherent
$\mathcal{O}_S$-algebra $\mathcal{B}$ and a quasi-coherent
$\mathcal{B}$-module $\mathcal{M}$ (as an
$\mathcal{O}_S$-module or as a $\mathcal{B}$-module, which comes to the
same thing (\hyperref[I.9.6.1]{I, 9.6.1})), there exists a pair
$(X,\mathcal{F})$ consisting of a prescheme $X$ affine over $S$ and a
quasi-coherent $\mathcal{O}_X$-module $\mathcal{F}$, such that
$\mathcal{A}(X)=\mathcal{B}$ and $\mathcal{A}(\mathcal{F})=\mathcal{M}$;
moreover, this pair is determined up to a unique isomorphism.
\end{proposition}

The uniqueness follows from (\hyperref[II.1.4.1]{1.4.1}) and
(\hyperref[II.1.4.2]{1.4.2}); the existence is proved as in
(\hyperref[II.1.3.1]{1.3.1}), and we again leave the details to the
reader. We denote by $\widetilde{\mathcal{M}}$ the $\mathcal{O}_X$-module
$\mathcal{F}$, and we say that it is \emph{associated} to the quasi-coherent
$\mathcal{B}$-module $\mathcal{M}$; for every affine open subset
$U\subset S$, $\widetilde{\mathcal{M}}|p^{-1}(U)$ (where $p$ is the
structural morphism $X\to S$) is canonically isomorphic to
$(\Gamma(U,\mathcal{M}))^{\sim}$.

\begin{corollary}[1.4.4]
\label{II.1.4.4}
In the category of quasi-coherent $\mathcal{B}$-modules,
$\widetilde{\mathcal{M}}$ is a covariant, additive and exact functor in
$\mathcal{M}$, which commutes with inductive limits and with direct sums.
\end{corollary}

Indeed, we are immediately reduced to the case where $S$ is affine, and the
corollary then follows from (\hyperref[I.1.3.5]{I, 1.3.5}),
(\hyperref[I.1.3.9]{I, 1.3.9}) and
(\hyperref[I.1.3.11]{I, 1.3.11}).

\begin{corollary}[1.4.5]
\label{II.1.4.5}
Under the hypotheses of (\hyperref[II.1.4.3]{1.4.3}), for
$\widetilde{\mathcal{M}}$ to be an $\mathcal{O}_X$-module of finite type, it
is necessary and sufficient that $\mathcal{M}$ be a $\mathcal{B}$-module of
finite type.
\end{corollary}

We are immediately reduced to the case where $S=\operatorname{Spec}(A)$ is
an affine scheme. Then $\mathcal{B}=\widetilde{B}$, where $B$ is an
$A$-algebra of finite type (\hyperref[I.9.6.2]{I, 9.6.2}), and
$\mathcal{M}=\widetilde{M}$, where $M$ is a $B$-module
(\hyperref[I.1.3.13]{I, 1.3.13}); \emph{on the prescheme $X$},
$\mathcal{O}_X$ is associated to the ring $B$ and
$\widetilde{\mathcal{M}}$ to the $B$-module $M$; for
$\widetilde{\mathcal{M}}$ to be of finite type, it is therefore necessary
and sufficient that $M$ be of finite type
(\hyperref[I.1.3.13]{I, 1.3.13}), whence our assertion.

\begin{proposition}[1.4.6]
\label{II.1.4.6}
Let $Y$ be a prescheme affine over $S$, and $X$, $X'$ two preschemes
affine over $Y$ (hence also over $S$
(\hyperref[II.1.3.5]{1.3.5})). Let
$\mathcal{B}=\mathcal{A}(Y)$, $\mathcal{A}=\mathcal{A}(X)$,
$\mathcal{A}'=\mathcal{A}(X')$. Then $X\times_Y X'$ is affine over
$Y$ (hence also over $S$) and $\mathcal{A}(X\times_Y X')$
is canonically identified with $\mathcal{A}\otimes_{\mathcal{B}}\mathcal{A}'$.
\end{proposition}

Indeed (\hyperref[I.9.6.1]{I, 9.6.1}),
$\mathcal{A}\otimes_{\mathcal{B}}\mathcal{A}'$ is a quasi-coherent
$\mathcal{B}$-algebra, hence also a quasi-coherent
$\mathcal{O}_S$-algebra
(\hyperref[I.9.6.1]{I, 9.6.1}); let $Z$ be the spectrum of
$\mathcal{A}\otimes_{\mathcal{B}}\mathcal{A}'$
(\hyperref[II.1.3.1]{1.3.1}). The canonical $\mathcal{B}$-homomorphisms
$\mathcal{A}\to\mathcal{A}\otimes_{\mathcal{B}}\mathcal{A}'$ and
$\mathcal{A}'\to\mathcal{A}\otimes_{\mathcal{B}}\mathcal{A}'$ correspond
(\hyperref[II.1.2.7]{1.2.7}) to $Y$-morphisms $p:Z\to X$ and
$p':Z\to X'$. To see that the triple $(Z,p,p')$ is a product
$X\times_Y X'$, we may restrict to the case where $S$ is an affine scheme
with ring $C$ (\hyperref[I.3.2.6.4]{I, 3.2.6.4}). But then $Y$, $X$,
$X'$ are affine schemes (\hyperref[II.1.3.4]{1.3.4}) whose
rings $B$, $A$, $A'$ are $C$-algebras such that
$\mathcal{B}=\widetilde{B}$, $\mathcal{A}=\widetilde{A}$,
$\mathcal{A}'=\widetilde{A'}$. We then know
(\hyperref[I.1.3.13]{I, 1.3.13}) that
$\mathcal{A}\otimes_{\mathcal{B}}\mathcal{A}'$ is canonically identified
with the $\mathcal{O}_S$-algebra $\widetilde{A\otimes_B A'}$, so the ring
$A(Z)$ is identified with $A\otimes_B A'$ and the morphisms $p$, $p'$
correspond to the canonical homomorphisms $A\to A\otimes_B A'$,
$A'\to A\otimes_B A'$. The proposition then follows from
(\hyperref[I.3.2.2]{I, 3.2.2}).

\begin{corollary}[1.4.7]
\label{II.1.4.7}
Let $\mathcal{F}$ (resp. $\mathcal{F}'$) be a quasi-coherent
$\mathcal{O}_X$-module (resp. $\mathcal{O}_{X'}$-module); then
$\mathcal{A}(\mathcal{F}\otimes_Y\mathcal{F}')$ is canonically identified
with $\mathcal{A}(\mathcal{F})\otimes_{\mathcal{A}(Y)}
\mathcal{A}(\mathcal{F}')$.
\end{corollary}

We know that $\mathcal{F}\otimes_Y\mathcal{F}'$ is quasi-coherent on
$X\times_Y X'$ (\hyperref[I.9.1.2]{I, 9.1.2}). Let $g:Y\to S$,
$f:X\to Y$, $f':X'\to Y$ be the structural morphisms, so that the
structural morphism
\oldpage[II]{11}
$h:Z\to S$ is equal to $g\circ f\circ p$ and to $g\circ f'\circ p'$. We
define a canonical homomorphism
\[
  \mathcal{A}(\mathcal{F})\otimes_{\mathcal{A}(Y)}\mathcal{A}(\mathcal{F}')
  \longrightarrow \mathcal{A}(\mathcal{F}\otimes_Y\mathcal{F}')
\]
as follows: for every open subset $U\subset S$, we have canonical
homomorphisms
\[
  \begin{aligned}
  \Gamma(f^{-1}(g^{-1}(U)),\mathcal{F})
  &\longrightarrow \Gamma(h^{-1}(U),p^*(\mathcal{F}))\quad\text{and}\quad
  \Gamma(f^{\prime-1}(g^{-1}(U)),\mathcal{F}')
  \longrightarrow \Gamma(h^{-1}(U),p^{\prime*}(\mathcal{F}')).
  \end{aligned}
\]
(\hyperref[0.4.4.3]{0, 4.4.3}), from which we deduce a canonical
homomorphism
\[
  \begin{aligned}
  &\Gamma(f^{-1}(g^{-1}(U)),\mathcal{F})
    \otimes_{\Gamma(g^{-1}(U),\mathcal{O}_Y)}
    \Gamma(f^{\prime-1}(g^{-1}(U)),\mathcal{F}')\longrightarrow\\
  &\hspace{17em}
    \Gamma(h^{-1}(U),p^*(\mathcal{F}))
    \otimes_{\Gamma(h^{-1}(U),\mathcal{O}_Z)}
    \Gamma(h^{-1}(U),p^{\prime*}(\mathcal{F}')).
  \end{aligned}
\]

To see that this indeed gives an isomorphism of $\mathcal{A}(Z)$-modules, we
may restrict to the case where $S$ (and consequently $X$, $X'$, $Y$,
$X\times_Y X'$) are affine schemes, and (with the notation of
(\hyperref[II.1.4.6]{1.4.6})), $\mathcal{F}=\widetilde{M}$,
$\mathcal{F}'=\widetilde{M'}$, $M$ (resp. $M'$) being an $A$-module (resp.
an $A'$-module). Then $\mathcal{F}\otimes_Y\mathcal{F}'$ is identified with
the sheaf on $X\times_Y X'$ associated to the $(A\otimes_B A')$-module
$M\otimes_B M'$ (\hyperref[I.9.1.3]{I, 9.1.3}), and the corollary follows
from the canonical identification of the $\mathcal{O}_S$-modules
$\widetilde{M\otimes_B M'}$ and
$\widetilde{M}\otimes_{\widetilde{B}}\widetilde{M'}$ (where $M$, $M'$ and
$B$ are regarded as $C$-modules)
(\hyperref[I.1.3.13]{I, 1.3.13} and
\hyperref[I.1.6.3]{I, 1.6.3}).

If we apply (\hyperref[II.1.4.7]{1.4.7}) in particular to the case where
$X=Y$ and $\mathcal{F}'=\mathcal{O}_{X'}$, we see that the
$\mathcal{A}'$-module $\mathcal{A}(f^{\prime*}(\mathcal{F}))$ is identified
with $\mathcal{A}(\mathcal{F})\otimes_{\mathcal{B}}\mathcal{A}'$.

\begin{env}[1.4.8]
\label{II.1.4.8}
In particular, when $X=X'=Y$ ($X$ being affine over $S$), we see
that if $\mathcal{F}$, $\mathcal{G}$ are two quasi-coherent
$\mathcal{O}_X$-modules, we have
\[
  \mathcal{A}(\mathcal{F}\otimes_{\mathcal{O}_X}\mathcal{G})
  =\mathcal{A}(\mathcal{F})\otimes_{\mathcal{A}(X)}\mathcal{A}(\mathcal{G})
  \tag{1.4.8.1}\label{II.1.4.8.1}
\]
up to canonical functorial isomorphism. If moreover $\mathcal{F}$ admits
a finite presentation, it follows from
(\hyperref[I.1.6.3]{I, 1.6.3} and
\hyperref[I.1.3.12]{I, 1.3.12}) that
\[
  \mathcal{A}(\mathcal{H}om_X(\mathcal{F},\mathcal{G}))
  =\mathcal{H}om_{\mathcal{A}(X)}
    (\mathcal{A}(\mathcal{F}),\mathcal{A}(\mathcal{G}))
  \tag{1.4.8.2}\label{II.1.4.8.2}
\]
up to canonical isomorphism.
\end{env}

\begin{remark}[1.4.9]
\label{II.1.4.9}
If $X$, $X'$ are two preschemes affine over $S$, the sum
$X\sqcup X'$ is also affine over $S$, the sum of two affine schemes
being an affine scheme.
\end{remark}

\begin{proposition}[1.4.10]
\label{II.1.4.10}
Let $S$ be a prescheme, $\mathcal{B}$ a quasi-coherent
$\mathcal{O}_S$-algebra, and $X=\operatorname{Spec}(\mathcal{B})$. For every
quasi-coherent sheaf of ideals $\mathcal{J}$ of $\mathcal{B}$,
$\widetilde{\mathcal{J}}$ is a quasi-coherent sheaf of ideals of
$\mathcal{O}_X$, and the closed sub-prescheme $Y$ of $X$ defined by
$\widetilde{\mathcal{J}}$ is canonically isomorphic to
$\operatorname{Spec}(\mathcal{B}/\mathcal{J})$.
\end{proposition}

Indeed, it follows immediately from (\hyperref[I.4.2.3]{I, 4.2.3}) that
$Y$ is affine over $S$; by virtue of
(\hyperref[II.1.3.1]{1.3.1}), we are therefore reduced to the case where
$S$ is affine, and the proposition then follows immediately from
(\hyperref[I.4.1.2]{I, 4.1.2}).

We can also express the result of
(\hyperref[II.1.4.10]{1.4.10}) by saying that if
$h:\mathcal{B}\to\mathcal{B}'$ is a \emph{surjective} homomorphism of
quasi-coherent $\mathcal{O}_S$-algebras, then
${}^a h:\operatorname{Spec}(\mathcal{B}')\to
\operatorname{Spec}(\mathcal{B})$ is a \emph{closed immersion}.

\oldpage[II]{12}
\begin{proposition}[1.4.11]
\label{II.1.4.11}
Let $S$ be a prescheme, $\mathcal{B}$ a quasi-coherent
$\mathcal{O}_S$-algebra, and $X=\operatorname{Spec}(\mathcal{B})$. For every
quasi-coherent ideal $\mathcal{K}$ of $\mathcal{O}_S$, we have (denoting by
$f$ the structural morphism $X\to S$)
$f^*(\mathcal{K})\mathcal{O}_X=\widetilde{\mathcal{K}\mathcal{B}}$ up to
canonical isomorphism.
\end{proposition}

The question being local on $S$, we are reduced to the case where
$S=\operatorname{Spec}(A)$ is affine, and in this case the proposition is
none other than (\hyperref[I.1.6.9]{I, 1.6.9}).

\subsection{Change of the base prescheme}
\label{II.1.5}

\begin{proposition}[1.5.1]
\label{II.1.5.1}
Let $X$ be a prescheme affine over $S$. For every extension
$g:S'\to S$ of the base prescheme, $X'=X_{(S')}=X\times_S S'$ is affine
over $S'$.
\end{proposition}

If $f'$ is the projection $X'\to S'$, it suffices to prove that
$f^{\prime-1}(U')$ is an affine open subset for every affine open subset
$U'$ of $S'$ such that $g(U')$ is contained in an affine open subset $U$ of
$S$ (\hyperref[II.1.2.1]{1.2.1}); we may therefore restrict to the case
where $S$ and $S'$ are affine, and it suffices to prove that $X'$ is then an
affine scheme (\hyperref[II.1.3.4]{1.3.4}). But then
(\hyperref[II.1.3.4]{1.3.4}) $X$ is an affine scheme, and if $A$, $A'$
and $B$ are the rings of $S$, $S'$ and $X$ respectively, we know that $X'$
is an affine scheme with ring $A'\otimes_A B$
(\hyperref[I.3.2.2]{I, 3.2.2}).

\begin{corollary}[1.5.2]
\label{II.1.5.2}
Under the hypotheses of (\hyperref[II.1.5.1]{1.5.1}), let $f:X\to S$
be the structural morphism, and $f':X'\to S'$, $g':X'\to X$ the projections,
so that the diagram
\[
  \xymatrix{
    X\ar[d] & X'\ar[l]_{g'}\ar[d]^{f'}\\
    S & S'\ar[l]_g
  }
\]
is commutative. For every quasi-coherent $\mathcal{O}_X$-module
$\mathcal{F}$, there exists a canonical isomorphism of
$\mathcal{O}_{S'}$-modules
\[
  u:g^*(f_*(\mathcal{F}))\xrightarrow{\sim}
  f'_*(g^{\prime*}(\mathcal{F})).
  \tag{1.5.2.1}\label{II.1.5.2.1}
\]
In particular, there exists a canonical isomorphism of $\mathcal{A}(X')$
onto $g^*(\mathcal{A}(X))$.
\end{corollary}

To define $u$, it suffices to define a homomorphism
\[
  v:f_*(\mathcal{F})\longrightarrow
  g_*(f'_*(g^{\prime*}(\mathcal{F})))
  =f_*(g'_*(g^{\prime*}(\mathcal{F})))
\]
and to take $u=v^\sharp$ (\hyperref[0.4.4.3]{0, 4.4.3}). We will take
$v=f_*(\rho)$, where $\rho$ is the canonical homomorphism
$\mathcal{F}\to g'_*(g^{\prime*}(\mathcal{F}))$
(\hyperref[0.4.4.3]{0, 4.4.3}). To prove that $u$ is an isomorphism,
we may restrict to the case where $S$ and $S'$, and hence $X$ and $X'$, are
affine; with the notation of (\hyperref[II.1.5.1]{1.5.1}), we then have
$\mathcal{F}=\widetilde{M}$, where $M$ is a $B$-module. We then see
immediately that $g^*(f_*(\mathcal{F}))$ and
$f'_*(g^{\prime*}(\mathcal{F}))$ are both equal to the
$\mathcal{O}_{S'}$-module associated to the $A'$-module $A'\otimes_A M$
(where $M$ is regarded as an $A$-module), and that $u$ is the homomorphism
associated to the identity (\hyperref[I.1.6.3]{I, 1.6.3},
\hyperref[I.1.6.5]{1.6.5} and \hyperref[I.1.6.7]{1.6.7}).

\begin{remark}[1.5.3]
\label{II.1.5.3}
One should beware of believing that (\hyperref[II.1.5.2]{1.5.2}) remains
valid when $X$ is no longer assumed to be affine over $S$, even when
$S'=\operatorname{Spec}(k(s))$ ($s\in S$) and $S'\to S$ is the
canonical morphism (\hyperref[I.2.4.5]{I, 2.4.5}) --- in which case $X'$
is none other than the \emph{fibre} $f^{-1}(s)$
(\hyperref[I.3.6.2]{I, 3.6.2}). In other words, when $X$ is not affine
over $S$, the operation
\oldpage[II]{13}
``direct image of quasi-coherent sheaves'' does not commute with the
operation of ``passage to fibres''. We will nevertheless see in Chapter~III
(\hyperref[III.4.2.4]{III, 4.2.4}) a result in this direction, of an
``asymptotic'' nature, valid for \emph{coherent} sheaves on $X$ when
$f$ is proper (\hyperref[II.5.4]{5.4}) and $S$ Noetherian.
\end{remark}

\begin{corollary}[1.5.4]
\label{II.1.5.4}
For every prescheme $X$ affine over $S$ and every $s\in S$, the fibre
$f^{-1}(s)$ (where $f$ denotes the structural morphism $X\to S$) is an
affine scheme.
\end{corollary}

It suffices to apply (\hyperref[II.1.5.1]{1.5.1}) to
$S'=\operatorname{Spec}(k(s))$ and to use
(\hyperref[II.1.3.4]{1.3.4}).

\begin{corollary}[1.5.5]
\label{II.1.5.5}
Let $X$ be an $S$-prescheme, and $S'$ a prescheme affine over $S$; then
$X'=X_{(S')}$ is a prescheme affine over $X$. Moreover, if
$f:X\to S$ is the structural morphism, there is a canonical isomorphism of
$\mathcal{O}_X$-algebras
$\mathcal{A}(X')\xrightarrow{\sim}f^*(\mathcal{A}(S'))$, and, for every
quasi-coherent $\mathcal{A}(S')$-module $\mathcal{M}$, a canonical
di-isomorphism
$f^*(\mathcal{M})\xrightarrow{\sim}
\mathcal{A}(f^{\prime*}(\widetilde{\mathcal{M}}))$, where we denote by
$f'=f_{(S')}$ the structural morphism $X'\to S'$.
\end{corollary}

It suffices to interchange the roles of $X$ and $S'$ in
(\hyperref[II.1.5.1]{1.5.1}) and
(\hyperref[II.1.5.2]{1.5.2}).

\begin{env}[1.5.6]
\label{II.1.5.6}
Now let $S$, $S'$ be two preschemes, $q:S'\to S$ a morphism,
$\mathcal{B}$ (resp. $\mathcal{B}'$) a quasi-coherent $\mathcal{O}_S$-algebra
(resp. $\mathcal{O}_{S'}$-algebra), and $u:\mathcal{B}\to\mathcal{B}'$
a $q$-morphism (that is, a homomorphism
$\mathcal{B}\to q_*(\mathcal{B}')$ of $\mathcal{O}_S$-algebras). If
$X=\operatorname{Spec}(\mathcal{B})$ and
$X'=\operatorname{Spec}(\mathcal{B}')$, we canonically deduce from this a
morphism
\[
  v=\operatorname{Spec}(u):X'\to X
\]
such that the diagram
\[
  \xymatrix{
    X'\ar[r]^v\ar[d] & X\ar[d]\\
    S'\ar[r]_q & S
  }
  \tag{1.5.6.1}\label{II.1.5.6.1}
\]
is commutative (the vertical arrows being the structural morphisms).
Indeed, giving $u$ is equivalent to giving a homomorphism of quasi-coherent
$\mathcal{O}_{S'}$-algebras
$u^\sharp:q^*(\mathcal{B})\to\mathcal{B}'$
(\hyperref[0.4.4.3]{0, 4.4.3}); it therefore canonically defines an
$S'$-morphism
\[
  w:\operatorname{Spec}(\mathcal{B}')\to
  \operatorname{Spec}(q^*(\mathcal{B}))
\]
such that $\mathcal{A}(w)=u^\sharp$
(\hyperref[II.1.2.7]{1.2.7}). On the other hand, it follows from
(\hyperref[II.1.5.2]{1.5.2}) that
$\operatorname{Spec}(q^*(\mathcal{B}))$ is canonically identified with
$X\times_S S'$; the morphism $v$ is the composite
$X'\xrightarrow{w}X\times_S S'\xrightarrow{p_1}X$ of $w$ with the first
projection, and the commutativity of
(\hyperref[II.1.5.6.1]{1.5.6.1}) follows from the definitions. Let $U$
(resp. $U'$) be an affine open subset of $S$ (resp. $S'$) such that
$q(U')\subset U$, $A=\Gamma(U,\mathcal{O}_S)$,
$A'=\Gamma(U',\mathcal{O}_{S'})$ their rings,
$B=\Gamma(U,\mathcal{B})$, $B'=\Gamma(U',\mathcal{B}')$; the restriction of
$u$ as a $(q|U')$-morphism
$\mathcal{B}|U\to\mathcal{B}'|U'$ corresponds to a di-homomorphism of
algebras $B\to B'$; if $V$, $V'$ are the inverse images of $U$, $U'$ in $X$,
$X'$ respectively under the structural morphisms, the morphism $V'\to V$
that is the restriction of $v$ corresponds (\hyperref[I.1.7.3]{I, 1.7.3})
to the preceding di-homomorphism.
\end{env}

\begin{env}[1.5.7]
\label{II.1.5.7}
Under the same hypotheses as in (\hyperref[II.1.5.6]{1.5.6}), let
$\mathcal{M}$ be a quasi-coherent $\mathcal{B}$-module; there is then a
canonical isomorphism of $\mathcal{O}_{X'}$-modules
\[
  v^*(\widetilde{\mathcal{M}})\xrightarrow{\sim}
  \bigl(q^*(\mathcal{M})\otimes_{q^*(\mathcal{B})}\mathcal{B}'\bigr)^{\sim}
  \tag{1.5.7.1}\label{II.1.5.7.1}
\]

\oldpage[II]{14}
Indeed, the canonical isomorphism
(\hyperref[II.1.5.2.1]{1.5.2.1}) provides a canonical isomorphism of
$p_1^*(\widetilde{\mathcal{M}})$ onto the sheaf on
$\operatorname{Spec}(q^*(\mathcal{B}))$ associated to the
$q^*(\mathcal{B})$-module $q^*(\mathcal{M})$, and it then suffices to apply
(\hyperref[II.1.4.7]{1.4.7}).
\end{env}

\subsection{Affine morphisms}
\label{II.1.6}

\begin{env}[1.6.1]
\label{II.1.6.1}
We will say that a morphism $f:X\to Y$ of preschemes is \emph{affine} if it
defines $X$ as a prescheme affine over $Y$. The properties of
preschemes affine over another one translate as follows into this
language:
\end{env}

\begin{proposition}[1.6.2]
\label{II.1.6.2}
\begin{enumerate}
  \item[\rm(i)] A closed immersion is affine.
  \item[\rm(ii)] The composite of two affine morphisms is affine.
  \item[\rm(iii)] If $f:X\to Y$ is an affine $S$-morphism,
    $f_{(S')}:X_{(S')}\to Y_{(S')}$ is affine for every extension $S'\to S$
    of the base.
  \item[\rm(iv)] If $f:X\to Y$ and $f':X'\to Y'$ are two affine
    $S$-morphisms,
    \[
      f\times_S f':X\times_S X'\to Y\times_S Y'
    \]
    is affine.
  \item[\rm(v)] If $f:X\to Y$ and $g:Y\to Z$ are two morphisms such that
    $g\circ f$ is affine and $g$ separated, then $f$ is affine.
  \item[\rm(vi)] If $f$ is affine, so is
    $f_{\mathrm{red}}$.
\end{enumerate}
\end{proposition}

By virtue of (\hyperref[I.5.5.12]{I, 5.5.12}), it suffices to prove (i),
(ii) and (iii). Now (i) is none other than Example
(\hyperref[II.1.2.2]{1.2.2}), and (ii) is none other than
(\hyperref[II.1.3.5]{1.3.5}); finally, (iii) follows from
(\hyperref[II.1.5.1]{1.5.1}) since $X_{(S')}$ is identified with the
product $X\times_Y Y_{(S')}$ (\hyperref[I.3.3.11]{I, 3.3.11}).

\begin{corollary}[1.6.3]
\label{II.1.6.3}
If $X$ is an affine scheme and $Y$ a scheme, every morphism $X\to Y$ is
affine.
\end{corollary}

\begin{proposition}[1.6.4]
\label{II.1.6.4}
Let $Y$ be a locally Noetherian prescheme, and $f:X\to Y$ a morphism of
finite type. For $f$ to be affine, it is necessary and sufficient that
$f_{\mathrm{red}}$ be so.
\end{proposition}

In view of (\hyperref[II.1.6.2]{1.6.2, (vi)}), we only have to prove the
sufficiency of the condition. It suffices to prove that if $Y$ is affine and
Noetherian, $X$ is affine; now $Y_{\mathrm{red}}$ is then affine, so the same
is true of $X_{\mathrm{red}}$, by hypothesis. Now, $X$ is Noetherian,
so the conclusion follows from (\hyperref[I.6.1.7]{I, 6.1.7}).


\subsection{Vector bundle associated to a sheaf of modules}
\label{II.1.7}

\begin{env}[1.7.1]
\label{II.1.7.1}
Let $A$ be a ring and $E$ an $A$-module. Recall that the \emph{symmetric
algebra} on $E$, denoted by $\mathbf{S}(E)$ (or $\mathbf{S}_A(E)$), is
the quotient algebra of the tensor algebra $\mathbf{T}(E)$ by the two-sided
ideal generated by the elements $x\otimes y-y\otimes x$, where $x,y$
run through $E$. The algebra $\mathbf{S}(E)$ is characterized by the
following universal property: if $\sigma$ is the canonical map
$E\to\mathbf{S}(E)$ (obtained by composing $E\to\mathbf{T}(E)$ with
the canonical map $\mathbf{T}(E)\to\mathbf{S}(E)$), every
$A$-linear map $E\to B$, where $B$ is a \emph{commutative} $A$-algebra,
factors uniquely as
$E\xrightarrow{\sigma}\mathbf{S}(E)\xrightarrow{g}B$, where $g$ is an
$A$-homomorphism \emph{of algebras}. We deduce immediately from this
characterization that for two $A$-modules $E$, $F$, we have
\[
  \mathbf{S}(E\oplus F)=\mathbf{S}(E)\otimes\mathbf{S}(F)
\]

\oldpage[II]{15}
up to canonical isomorphism; moreover, $\mathbf{S}(E)$ is a covariant
functor in $E$, from the category of $A$-modules to that of commutative
$A$-algebras; finally, the preceding characterization also shows that if
$E=\varinjlim E_\lambda$, we have
$\mathbf{S}(E)=\varinjlim\mathbf{S}(E_\lambda)$ up to canonical
isomorphism. By abuse of language, a product
$\sigma(x_1)\sigma(x_2)\cdots\sigma(x_n)$, where the $x_i\in E$, is often
written $x_1x_2\cdots x_n$ if no confusion results. The algebra
$\mathbf{S}(E)$ is \emph{graded}, $\mathbf{S}_n(E)$ being the set of
linear combinations of products of $n$ elements of $E$ ($n\geq 0$);
the algebra $\mathbf{S}(A)$ is canonically isomorphic to the polynomial
algebra $A[T]$ in one indeterminate, and the algebra $\mathbf{S}(A^n)$ to
the polynomial algebra in $n$ indeterminates $A[T_1,\ldots,T_n]$.
\end{env}

\begin{env}[1.7.2]
\label{II.1.7.2}
Let $\varphi$ be a ring homomorphism $A\to B$. If $F$ is a $B$-module,
the canonical map $F\to\mathbf{S}(F)$ gives a canonical map
$F_{[\varphi]}\to\mathbf{S}(F)_{[\varphi]}$, which therefore factors as
$F_{[\varphi]}\to\mathbf{S}(F_{[\varphi]})\to
\mathbf{S}(F)_{[\varphi]}$; the canonical homomorphism
$\mathbf{S}(F_{[\varphi]})\to\mathbf{S}(F)_{[\varphi]}$ is surjective, but
not necessarily bijective. If $E$ is an $A$-module, every di-homomorphism
$E\to F$ (that is, every $A$-homomorphism $E\to F_{[\varphi]}$) therefore
canonically gives an $A$-homomorphism of algebras
$\mathbf{S}(E)\to\mathbf{S}(F_{[\varphi]})\to
\mathbf{S}(F)_{[\varphi]}$, that is, a di-homomorphism of algebras
$\mathbf{S}(E)\to\mathbf{S}(F)$.

With the same notation, for every $A$-module $E$,
$\mathbf{S}(E\otimes_A B)$ is canonically identified with the algebra
$\mathbf{S}(E)\otimes_A B$; this follows immediately from the universal
property of $\mathbf{S}(E)$ (\hyperref[II.1.7.1]{1.7.1}).
\end{env}

\begin{env}[1.7.3]
\label{II.1.7.3}
Let $R$ be a multiplicative subset of the ring $A$; applying
(\hyperref[II.1.7.2]{1.7.2}) to the ring $B=R^{-1}A$, and recalling that
$R^{-1}E=E\otimes_A R^{-1}A$, we see that we have
$\mathbf{S}(R^{-1}E)=R^{-1}\mathbf{S}(E)$ up to canonical isomorphism.
Moreover, if $R'\supset R$ is a second multiplicative subset of $A$, the
diagram
\[
  \xymatrix{
    R^{-1}E\ar[r]\ar[d] & {R'}^{-1}E\ar[d]\\
    \mathbf{S}(R^{-1}E)\ar[r] & \mathbf{S}({R'}^{-1}E)
  }
\]
is commutative.
\end{env}

\begin{env}[1.7.4]
\label{II.1.7.4}
Now let $(S,\mathcal{A})$ be a ringed space, and let $\mathcal{E}$ be an
$\mathcal{A}$-module on $S$. If to every open subset $U\subset S$ we
associate the $\Gamma(U,\mathcal{A})$-module
$\mathbf{S}(\Gamma(U,\mathcal{E}))$, we define (in view of the functorial
character of $\mathbf{S}(E)$ (\hyperref[II.1.7.2]{1.7.2})) a presheaf of
algebras; we say that the associated sheaf, which we denote by
$\mathbf{S}(\mathcal{E})$ or $\mathbf{S}_{\mathcal{A}}(\mathcal{E})$, is
the \emph{symmetric} $\mathcal{A}$-algebra of the $\mathcal{A}$-module
$\mathcal{E}$. It follows immediately from (\hyperref[II.1.7.1]{1.7.1})
that $\mathbf{S}(\mathcal{E})$ is the solution of a universal problem:
every homomorphism of $\mathcal{A}$-modules
$\mathcal{E}\to\mathcal{B}$, where $\mathcal{B}$ is a commutative
$\mathcal{A}$-algebra, factors in exactly one way as
$\mathcal{E}\to\mathbf{S}(\mathcal{E})\to\mathcal{B}$, the second arrow
being a homomorphism of $\mathcal{A}$-algebras. There is therefore a
bijective correspondence between homomorphisms $\mathcal{E}\to\mathcal{B}$
of $\mathcal{A}$-modules and homomorphisms
$\mathbf{S}(\mathcal{E})\to\mathcal{B}$ of $\mathcal{A}$-algebras. In
particular, every homomorphism $u:\mathcal{E}\to\mathcal{F}$ of
$\mathcal{A}$-modules defines a homomorphism
$\mathbf{S}(u):\mathbf{S}(\mathcal{E})\to\mathbf{S}(\mathcal{F})$ of
$\mathcal{A}$-algebras, and $\mathbf{S}(\mathcal{E})$ is therefore a
covariant functor in $\mathcal{E}$.

\oldpage[II]{16}
By virtue of (\hyperref[II.1.7.2]{1.7.2}) and of the commutation of
$\mathbf{S}$ with inductive limits, we have
$(\mathbf{S}(\mathcal{E}))_s=\mathbf{S}(\mathcal{E}_s)$ for every point
$s\in S$. If $\mathcal{E}$, $\mathcal{F}$ are two $\mathcal{A}$-modules,
$\mathbf{S}(\mathcal{E}\oplus\mathcal{F})$ is canonically identified with
$\mathbf{S}(\mathcal{E})\otimes_{\mathcal{A}}\mathbf{S}(\mathcal{F})$,
as one sees on the corresponding presheaves.

Note also that $\mathbf{S}(\mathcal{E})$ is a graded $\mathcal{A}$-algebra,
the infinite direct sum of the $\mathbf{S}_n(\mathcal{E})$, where the
$\mathcal{A}$-module $\mathbf{S}_n(\mathcal{E})$ is the sheaf associated to
the presheaf $U\mapsto\mathbf{S}_n(\Gamma(U,\mathcal{E}))$. If in
particular we take $\mathcal{E}=\mathcal{A}$, we see that
$\mathbf{S}_{\mathcal{A}}(\mathcal{A})$ is identified with
$\mathcal{A}[T]=\mathcal{A}\otimes_{\mathbf{Z}}\mathbf{Z}[T]$
($T$ an indeterminate, $\mathbf{Z}$ being regarded as a simple sheaf).
\end{env}

\begin{env}[1.7.5]
\label{II.1.7.5}
Let $(T,\mathcal{B})$ be a second ringed space and $f$ a morphism
$(S,\mathcal{A})\to(T,\mathcal{B})$. If $\mathcal{F}$ is a
$\mathcal{B}$-module, $\mathbf{S}(f^*(\mathcal{F}))$ is canonically
identified with $f^*(\mathbf{S}(\mathcal{F}))$; indeed, if
$f=(\psi,\theta)$, we have by definition
(\hyperref[0.4.3.1]{0, 4.3.1})
\[
  \mathbf{S}(f^*(\mathcal{F}))
  =\mathbf{S}(\psi^*(\mathcal{F})
    \otimes_{\psi^*(\mathcal{B})}\mathcal{A})
  =\mathbf{S}(\psi^*(\mathcal{F}))
    \otimes_{\psi^*(\mathcal{B})}\mathcal{A}
\]
(\hyperref[II.1.7.2]{1.7.2}); for every open subset $U$ of $S$ and
every section $h$ of $\mathbf{S}(\psi^*(\mathcal{F}))$ over $U$, $h$
coincides, in a neighbourhood $V$ of each point $s\in U$, with an element of
$\mathbf{S}(\Gamma(V,\psi^*(\mathcal{F})))$; if we refer to the
definition of $\psi^*(\mathcal{F})$
(\hyperref[0.3.7.1]{0, 3.7.1}) and take into account that every element
of $\mathbf{S}(E)$, for a module $E$, is a linear combination of finitely
many products of elements of $E$, we see that there are a neighbourhood $W$
of $\psi(s)$ in $T$, a section $h'$ of $\mathbf{S}(\mathcal{F})$ over
$W$ and a neighbourhood $V'\subset V\cap\psi^{-1}(W)$ of $s$ such that $h$
coincides with $t\mapsto h'(\psi(t))$ over $V'$; whence our
assertion.
\end{env}

\begin{proposition}[1.7.6]
\label{II.1.7.6}
Let $A$ be a ring, $S=\operatorname{Spec}(A)$ its prime spectrum, and
$\mathcal{E}=\widetilde{M}$ the $\mathcal{O}_S$-module associated to an
$A$-module $M$; then the $\mathcal{O}_S$-algebra $\mathbf{S}(\mathcal{E})$ is
associated to the $A$-algebra $\mathbf{S}(M)$.
\end{proposition}

\begin{proof}
Indeed, for every $f\in A$,
$\mathbf{S}(M_f)=(\mathbf{S}(M))_f$
(\hyperref[II.1.7.3]{1.7.3}), and the proposition therefore follows from
the definitions (\hyperref[I.1.3.4]{I, 1.3.4}).
\end{proof}

\begin{corollary}[1.7.7]
\label{II.1.7.7}
If $S$ is a prescheme and $\mathcal{E}$ a quasi-coherent
$\mathcal{O}_S$-module, the $\mathcal{O}_S$-algebra $\mathbf{S}(\mathcal{E})$
is quasi-coherent. If moreover $\mathcal{E}$ is of finite type, each of the
$\mathcal{O}_S$-modules $\mathbf{S}_n(\mathcal{E})$ is of finite type.
\end{corollary}

\begin{proof}
The first assertion is an immediate consequence of
(\hyperref[II.1.7.6]{1.7.6}) and
(\hyperref[I.1.4.1]{I, 1.4.1}); the second follows from the fact that,
if $E$ is an $A$-module of finite type, $\mathbf{S}_n(E)$ is an $A$-module
of finite type; we then apply (\hyperref[I.1.3.13]{I, 1.3.13}).
\end{proof}

\begin{definition}[1.7.8]
\label{II.1.7.8}
Let $\mathcal{E}$ be a quasi-coherent $\mathcal{O}_S$-module. We call
\emph{vector bundle over $S$ defined by $\mathcal{E}$}, and denote by
$\mathbf{V}(\mathcal{E})$, the spectrum
(\hyperref[II.1.3.1]{1.3.1}) of the quasi-coherent $\mathcal{O}_S$-algebra
$\mathbf{S}(\mathcal{E})$.
\end{definition}

By virtue of (\hyperref[II.1.2.7]{1.2.7}), for every $S$-prescheme $X$,
there is a canonical bijective correspondence between the $S$-morphisms
$X\to\mathbf{V}(\mathcal{E})$ and the homomorphisms of
$\mathcal{O}_S$-algebras $\mathbf{S}(\mathcal{E})\to\mathcal{A}(X)$, hence
also between these $S$-morphisms and the homomorphisms of
$\mathcal{O}_S$-modules
$\mathcal{E}\to\mathcal{A}(X)=f_*(\mathcal{O}_X)$ (where $f$ is the
structural morphism $X\to S$). In particular:

\begin{env}[1.7.9]
\label{II.1.7.9}
Take for $X$ a sub-prescheme induced by $S$ on an open subset
$U\subset S$. Then the $S$-morphisms $U\to\mathbf{V}(\mathcal{E})$ are none
other than the $U$-sections (\hyperref[I.2.5.5]{I, 2.5.5}) of the
$U$-prescheme induced by $\mathbf{V}(\mathcal{E})$ on the open subset
$p^{-1}(U)$ (where $p$ is the structural morphism
$\mathbf{V}(\mathcal{E})\to S$). By what we have just seen, these
$U$-sections correspond bijectively to the homomorphisms of
$\mathcal{O}_S$-modules
$\mathcal{E}\to j_*(\mathcal{O}_S|U)$ (where $j$ is the canonical injection
$U\to S$), or,

\oldpage[II]{17}
what amounts to the same (\hyperref[0.4.4.3]{0, 4.4.3}), to the
$(\mathcal{O}_S|U)$-homomorphisms
$j^*(\mathcal{E})=\mathcal{E}|U\to\mathcal{O}_S|U$. Moreover, it is
immediate that to the restriction to an open subset $U'\subset U$ of an
$S$-morphism $U\to\mathbf{V}(\mathcal{E})$ there corresponds the
restriction to $U'$ of the corresponding homomorphism
$\mathcal{E}|U\to\mathcal{O}_S|U$. We conclude that \emph{the sheaf of
germs of $S$-sections} of $\mathbf{V}(\mathcal{E})$ is canonically
identified with the \emph{dual $\check{\mathcal{E}}$} of $\mathcal{E}$.

In particular, if we take $X=U=S$, to the \emph{zero} homomorphism
$\mathcal{E}\to\mathcal{O}_S$ there corresponds a canonical $S$-section of
$\mathbf{V}(\mathcal{E})$, called the \emph{zero $S$-section}
(cf. (\hyperref[II.8.3.3]{8.3.3})).
\end{env}

\begin{env}[1.7.10]
\label{II.1.7.10}
Now take for $X$ the spectrum $\{\xi\}$ of a field $K$; the
structural morphism $f:X\to S$ therefore corresponds to a monomorphism
$k(s)\to K$, where $s=f(\xi)$
(\hyperref[I.2.4.6]{I, 2.4.6}); the $S$-morphisms
$\{\xi\}\to\mathbf{V}(\mathcal{E})$ are therefore none other than the
\emph{geometric points of $\mathbf{V}(\mathcal{E})$ with values in
the extension $K$ of $k(s)$} (\hyperref[I.3.4.5]{I, 3.4.5}), points
which are localized at the points of $p^{-1}(s)$. The set of these points,
which we may also call \emph{the geometric fibre rational over $K$ of
$\mathbf{V}(\mathcal{E})$ above the point $s$}, is identified by
(\hyperref[II.1.7.8]{1.7.8}) with the set of homomorphisms of
$\mathcal{O}_S$-modules
$\mathcal{E}\to f_*(\mathcal{O}_X)$, or, what amounts to the same
(\hyperref[0.4.4.3]{0, 4.4.3}), with the set of homomorphisms of
$\mathcal{O}_X$-modules
$f^*(\mathcal{E})\to\mathcal{O}_X=K$. But by definition
(\hyperref[0.4.3.1]{0, 4.3.1}) we have
$f^*(\mathcal{E})=\mathcal{E}_s\otimes_{\mathcal{O}_s}K
=\mathcal{E}^s\otimes_{k(s)}K$, setting
$\mathcal{E}^s=\mathcal{E}_s/\mathfrak{m}_s\mathcal{E}_s$; the geometric
fibre of $\mathbf{V}(\mathcal{E})$ rational over $K$ above
$s$ is therefore identified with the \emph{dual} of the \emph{$K$-vector
space $\mathcal{E}^s\otimes_{k(s)}K$}; if $\mathcal{E}^s$ or $K$ is of
finite dimension over $k(s)$, this dual is also identified with
$(\mathcal{E}^s)^\vee\otimes_{k(s)}K$, where
$(\mathcal{E}^s)^\vee$ denotes the dual of the $k(s)$-vector space
$\mathcal{E}^s$.
\end{env}

\begin{proposition}[1.7.11]
\label{II.1.7.11}
\begin{enumerate}
  \item[\rm(i)] $\mathbf{V}(\mathcal{E})$ is a contravariant functor in
    $\mathcal{E}$ from the category of quasi-coherent
    $\mathcal{O}_S$-modules to the category of affine $S$-schemes.
  \item[\rm(ii)] If $\mathcal{E}$ is an $\mathcal{O}_S$-module of finite
    type, $\mathbf{V}(\mathcal{E})$ is of finite type over $S$.
  \item[\rm(iii)] If $\mathcal{E}$ and $\mathcal{F}$ are two
    quasi-coherent $\mathcal{O}_S$-modules,
    $\mathbf{V}(\mathcal{E}\oplus\mathcal{F})$ is canonically identified
    with $\mathbf{V}(\mathcal{E})\times_S\mathbf{V}(\mathcal{F})$.
  \item[\rm(iv)] Let $g:S'\to S$ be a morphism; for every
    quasi-coherent $\mathcal{O}_S$-module $\mathcal{E}$,
    $\mathbf{V}(g^*(\mathcal{E}))$ is canonically identified with
    $\mathbf{V}(\mathcal{E})_{(S')}
    =\mathbf{V}(\mathcal{E})\times_S S'$.
  \item[\rm(v)] To a surjective homomorphism
    $\mathcal{E}\to\mathcal{F}$ of quasi-coherent $\mathcal{O}_S$-modules
    there corresponds a closed immersion
    $\mathbf{V}(\mathcal{F})\to\mathbf{V}(\mathcal{E})$.
\end{enumerate}
\end{proposition}

\begin{proof}
(i) is an immediate consequence of
(\hyperref[II.1.2.7]{1.2.7}), taking into account that every
homomorphism of $\mathcal{O}_S$-modules $\mathcal{E}\to\mathcal{F}$
canonically defines a homomorphism of $\mathcal{O}_S$-algebras
$\mathbf{S}(\mathcal{E})\to\mathbf{S}(\mathcal{F})$. (ii) follows
immediately from the definitions (\hyperref[I.6.3.1]{I, 6.3.1}) and
from the fact that if $E$ is an $A$-module of finite type, $\mathbf{S}(E)$
is an $A$-algebra of finite type. To prove (iii), it suffices to start
from the canonical isomorphism
$\mathbf{S}(\mathcal{E}\oplus\mathcal{F})
\simeq\mathbf{S}(\mathcal{E})
\otimes_{\mathcal{O}_S}\mathbf{S}(\mathcal{F})$
(\hyperref[II.1.7.4]{1.7.4}) and to apply
(\hyperref[II.1.4.6]{1.4.6}). Similarly, to prove (iv), it suffices to
start from the canonical isomorphism
$\mathbf{S}(g^*(\mathcal{E}))
\simeq g^*(\mathbf{S}(\mathcal{E}))$
(\hyperref[II.1.7.5]{1.7.5}) and to apply
(\hyperref[II.1.5.2]{1.5.2}). Finally, to establish (v), it suffices to
remark that if the homomorphism $\mathcal{E}\to\mathcal{F}$ is surjective,
so is the corresponding homomorphism
$\mathbf{S}(\mathcal{E})\to\mathbf{S}(\mathcal{F})$ of
$\mathcal{O}_S$-algebras, and the conclusion follows from
(\hyperref[II.1.4.10]{1.4.10}).
\end{proof}

\begin{env}[1.7.12]
\label{II.1.7.12}
Take in particular $\mathcal{E}=\mathcal{O}_S$; the prescheme
$\mathbf{V}(\mathcal{O}_S)$ is the affine $S$-scheme which is the spectrum
of the $\mathcal{O}_S$-algebra $\mathbf{S}(\mathcal{O}_S)$, and the latter
is identified with the $\mathcal{O}_S$-algebra
$\mathcal{O}_S[T]=\mathcal{O}_S\otimes_{\mathbf{Z}}\mathbf{Z}[T]$
\oldpage[II]{18}
($T$ an indeterminate); this is evident when
$S=\operatorname{Spec}(\mathbf{Z})$, by virtue of
(\hyperref[II.1.7.6]{1.7.6}), and we pass from there to the general case
by considering the structural morphism
$S\to\operatorname{Spec}(\mathbf{Z})$ and using
(\hyperref[II.1.7.11]{1.7.11, (iv)}). Because of this result, we also
set $\mathbf{V}(\mathcal{O}_S)=S[T]$, and we therefore have the formula
\[
  S[T]=S\otimes_{\mathbf{Z}}\mathbf{Z}[T].
  \tag{1.7.12.1}\label{II.1.7.12.1}
\]

The identification of the sheaf of germs of $S$-sections of $S[T]$ with
$\mathcal{O}_S$, already seen in
(\hyperref[I.3.3.15]{I, 3.3.15}), is recovered here in a more
general context, as a particular case of (\hyperref[II.1.7.9]{1.7.9}).
\end{env}

\begin{env}[1.7.13]
\label{II.1.7.13}
For every $S$-prescheme $X$, we have seen
(\hyperref[II.1.7.8]{1.7.8}) that
$\operatorname{Hom}_S(X,S[T])$ is canonically identified with
$\operatorname{Hom}_{\mathcal{O}_S}(\mathcal{O}_S,\mathcal{A}(X))$, which
is canonically isomorphic to $\Gamma(S,\mathcal{A}(X))$, and is
consequently endowed with a ring structure; moreover, to every
$S$-morphism $h:X\to Y$ there corresponds a morphism
$\Gamma(\mathcal{A}(h)):\Gamma(S,\mathcal{A}(Y))
\to\Gamma(S,\mathcal{A}(X))$ for the ring structures
(\hyperref[II.1.1.2]{1.1.2}). When we endow
$\operatorname{Hom}_S(X,S[T])$ with the ring structure thus defined, we
therefore see that we can regard $\operatorname{Hom}_S(X,S[T])$ as a
\emph{contravariant functor in $X$}, from the category of $S$-preschemes
to that of rings. On the other hand,
$\operatorname{Hom}_S(X,\mathbf{V}(\mathcal{E}))$ is likewise identified
with $\operatorname{Hom}_{\mathcal{O}_S}(\mathcal{E},\mathcal{A}(X))$
(where $\mathcal{A}(X)$ is regarded as an
\emph{$\mathcal{O}_S$-module}); we can consequently endow it canonically
with a structure of \emph{module} over the ring
$\operatorname{Hom}_S(X,S[T])$, and we see as above that the pair
\[
  \bigl(\operatorname{Hom}_S(X,S[T]),
  \operatorname{Hom}_S(X,\mathbf{V}(\mathcal{E}))\bigr)
\]
is a contravariant functor in $X$, with values in the category whose
elements are the pairs $(A,M)$ consisting of a ring $A$ and an $A$-module
$M$, the morphisms being the di-homomorphisms.

We will interpret these facts by saying that $S[T]$ is an
\emph{$S$-scheme of rings} and that $\mathbf{V}(\mathcal{E})$ is an
\emph{$S$-scheme of modules} over the $S$-scheme of rings $S[T]$
(cf. chap.~0, \textsection8).
\end{env}

\begin{env}[1.7.14]
\label{II.1.7.14}
We will see that the structure of $S$-scheme of modules defined on the
$S$-scheme $\mathbf{V}(\mathcal{E})$ makes it possible to reconstruct the
$\mathcal{O}_S$-module $\mathcal{E}$ up to unique isomorphism: to this
end, we will show that $\mathcal{E}$ is canonically isomorphic to a
sub-$\mathcal{O}_S$-module of
$\mathbf{S}(\mathcal{E})=\mathcal{A}(\mathbf{V}(\mathcal{E}))$ defined by
means of this structure. Indeed
(\hyperref[II.1.7.4]{1.7.4}), the set
$\operatorname{Hom}_{\mathcal{O}_S}(\mathbf{S}(\mathcal{E}),
\mathcal{A}(X))$ of homomorphisms of
\emph{$\mathcal{O}_S$-algebras} is canonically identified with
$\operatorname{Hom}_{\mathcal{O}_S}(\mathcal{E},\mathcal{A}(X))$, a set
of homomorphisms of \emph{$\mathcal{O}_S$-modules}: if $h$, $h'$ are two
elements of the latter set, $s_i$ ($1\leq i\leq n$) sections of
$\mathcal{E}$ over an open subset $U\subset S$, and $t$ a section of
$\mathcal{A}(X)$ over $U$, we have by definition
\[
  (h+h')(s_1s_2\cdots s_n)
  =\prod_{i=1}^n\bigl(h(s_i)+h'(s_i)\bigr)
\]
and
\[
  (t.h)(s_1s_2\cdots s_n)=t^n\prod_{i=1}^n h(s_i).
\]

This being so, if $z$ is a section of $\mathbf{S}(\mathcal{E})$ over
$U$, $h\mapsto h(z)$ is a map from
$\operatorname{Hom}_S(X,\mathbf{V}(\mathcal{E}))
=\operatorname{Hom}_{\mathcal{O}_S}(\mathbf{S}(\mathcal{E}),
\mathcal{A}(X))$ to $\Gamma(U,\mathcal{A}(X))$. We will
\oldpage[II]{19}
show that $\mathcal{E}$ is identified with the submodule of
$\mathbf{S}(\mathcal{E})$ such that, \emph{for every open subset
$U\subset S$, every section $z$ of this sub-$\mathcal{O}_S$-module over $U$
and every $S$-prescheme $X$, the map $h\mapsto h(z)$ from
$\operatorname{Hom}_{\mathcal{O}_S|U}(\mathbf{S}(\mathcal{E})|U,
\mathcal{A}(X)|U)$ to $\Gamma(U,\mathcal{A}(X))$ is a homomorphism of
$\Gamma(U,\mathcal{A}(X))$-modules}.

It is indeed immediate that $\mathcal{E}$ has this property; to prove
the converse, we may restrict to proving that when
$S=\operatorname{Spec}(A)$, $\mathcal{E}=\widetilde{M}$, a section $z$ of
$\mathbf{S}(\mathcal{E})$ over $S$ which (for $U=S$) has the property
stated above is necessarily a section of
$\mathcal{E}$; we then have $z=\sum_{n=0}^{\infty}z_n$, where
$z_n\in\mathbf{S}_n(M)$, and we must prove that $z_n=0$ for $n\ne1$.
Set $B=\mathbf{S}(M)$ and take for $X$ the prescheme
$\operatorname{Spec}(B[T])$, where $T$ is an indeterminate. The set
$\operatorname{Hom}_{\mathcal{O}_S}(\mathbf{S}(\mathcal{E}),
\mathcal{A}(X))$ is identified with the set of ring homomorphisms
$h:B\to B[T]$ (\hyperref[I.1.3.13]{I, 1.3.13}), and by what we saw
above, we have
$(T.h)(z)=\sum_{n=0}^{\infty}T^nh(z_n)$: the hypothesis on $z$ implies
that
$\sum_{n=0}^{\infty}T^nh(z_n)=T.\sum_{n=0}^{\infty}h(z_n)$ for every
homomorphism $h$. Taking in particular for $h$ the canonical injection, we
get $\sum_{n=0}^{\infty}T^nz_n=T.\sum_{n=0}^{\infty}z_n$, which indeed
implies the conclusion $z_n=0$ for $n\ne1$.
\end{env}

\begin{proposition}[1.7.15]
\label{II.1.7.15}
Let $Y$ be a prescheme whose underlying space is Noetherian, or a
quasi-compact scheme. Every affine $Y$-scheme $X$ of finite type over $Y$ is
$Y$-isomorphic to a closed sub-$Y$-scheme of a $Y$-scheme of the form
$\mathbf{V}(\mathcal{E})$, where $\mathcal{E}$ is a quasi-coherent
$\mathcal{O}_Y$-module of finite type.
\end{proposition}

\begin{proof}
Indeed, the quasi-coherent $\mathcal{O}_Y$-algebra $\mathcal{A}(X)$ is of
finite type (\hyperref[II.1.3.7]{1.3.7}). The hypotheses made imply
that $\mathcal{A}(X)$ is generated by a quasi-coherent
sub-$\mathcal{O}_Y$-module of finite type $\mathcal{E}$
(\hyperref[I.9.6.5]{I, 9.6.5}); by definition, this implies that
the canonical homomorphism
$\mathbf{S}(\mathcal{E})\to\mathcal{A}(X)$ canonically extending
the injection $\mathcal{E}\to\mathcal{A}(X)$ is \emph{surjective}; the
conclusion then follows from (\hyperref[II.1.4.10]{1.4.10}).
\end{proof}

